\cotefile{157-5}
\scanpage{389}\mspage{303}

% slanted
\margin{Refondre avec 4.13.23 etc.}

\uline{Définition} 4.13.11. Soient $M$ une catégorie \deleted{$\pi_{\unsure{0}}$}-accessible, $\pi_0$ un cardinal utile pour $M$, $\pi\inserted{\unsure{>|\geq}\pi_0}$ un cardinal. \deleted{\unsure{infini}} Une catégorie $I$ est dite $\pi_0$-$\pi$-\uline{adaptée} si a) $I$ est équivalente à une petite catégorie, b) $I$ est \unsure{grde|grd} devant $\pi_0$, et c) $I$ est stable par $\underset{\longrightarrow}{\lim}$ de cardinal $\leq\pi$, et \unsure{grdes|grde} devant $\pi_0$.

% slanted
\margin{\uline{NB} Si $\pi=\pi_0$, la condition c) équivaut à dire que $I$ est stable par facteurs directs.}

Un système inductif $\inserted{(x_i)_{i\in I}}$ de type $I$ dans $M$ est dit $\pi_0$-$\pi$-adapté, si $I$ est $\pi_0$-$\pi$-adapté i.e. satisfait les conditions a) b) c), et si de plus on a d) le foncteur $i\longrightarrow x_i:I\longrightarrow M$ commute aux $\underset{\longrightarrow}{\lim}$ (\unsure{grdes|grde} devant $\pi_0$ et de cardinal $\leq\pi$) envisagées dans c), et e) les $x_i$ sont dans $M_\pi$.

\uline{Proposition} 4.13.12. Soient $M$ une catégorie $\pi_0$-accessible, et $\pi>\pi_0$ un cardinal utile pour $M$. Alors tout $x\in\mathrm{Ob}M$ admet une représentation comme $\deleted{x\simeq}\underset{j\in J}{\underset{\longrightarrow}{\lim}}\,x_j$, où le système inductif $(x_j)_{j\in J}$ dans $M$ est $\pi_0$-$\pi$-adapté. De façon précise, on peut prendre $\deleted{\unsure{I}}J=M_\pi/x$.

% archive sheet number: 388
